下面的核函数把四个整数复制到输出缓冲区，主机程序逐个比较输入与输出。将完整代码保存为 \co@verify.co@：

\begin{lstlisting}[language=croqtile,caption={验证安装的最小程序}]
__co__ s32 [4] identity(s32 [4] input) {
  s32 [input.span] output;
  parallel i by 4
    output.at(i) = input.at(i);
  return output;
}
int main() {
  auto input = croq::make_spandata<croq::s32>(4);
  input[0] = 1; input[1] = 2; input[2] = 3; input[3] = 4;
  auto result = identity(input.view());
  for (int i = 0; i < 4; ++i)
    if (input[i] != result[i]) { std::cerr << "FAIL\n"; return 1; }
  std::cout << "OK\n";
}
\end{lstlisting}

在文件所在目录中编译并运行：

\begin{lstlisting}[language=croqtile,caption={编译并运行验证程序}]
croq verify.co -o verify
./verify
\end{lstlisting}

如果终端输出 \co@OK@，说明这次验证已完成编译、链接和运行，结果也符合预期。1.3 节介绍如何选择输出文件、目标平台和其他编译选项。
